% Generated draft: requires mathematical and source review.
\documentclass[11pt]{article}
\usepackage[margin=25mm]{geometry}
\usepackage{fontspec}
\setmainfont{Latin Modern Roman}
\usepackage{amsmath,amssymb,mathtools,unicode-math}
\setmathfont{Latin Modern Math}
\usepackage{enumitem,needspace,xurl,hyperref}
\setlength{\parindent}{0pt}
\setlength{\parskip}{5pt}
\setlength{\emergencystretch}{3em}
\newcommand{\N}{\mathbb{N}}
\newcommand{\Z}{\mathbb{Z}}
\newcommand{\Q}{\mathbb{Q}}
\newcommand{\R}{\mathbb{R}}
\newcommand{\C}{\mathbb{C}}
\newcommand{\F}{\mathbb{F}}
\newcommand{\A}{\mathbb{A}}
\newcommand{\PP}{\mathbb{P}}
\newcommand{\E}{\mathbb{E}}
\newcommand{\eps}{\varepsilon}
\newcommand{\dd}{\,\mathrm{d}}
\begin{document}
\section*{UnsolvedMath ID: 11000101}
\textbf{Source problem code:} catalogue entry
\subsection{Definitions and mathematical statement}
Let \(S\) be the surface of the source and \(\operatorname{Mod}_S\) its mapping class group; the genus, punctures, boundary and orientation conventions are not supplied. A discrete group \(G\) has Kazhdan's property \((T)\) if every unitary representation \(\pi:G\to\mathrm U(\mathcal H)\) with almost invariant unit vectors has a nonzero invariant vector. Here almost invariance means
\[
 \forall Q\subset G\text{ finite}\ \forall\eps>0\ \exists v\in\mathcal H:
 \quad\|v\|=1,\quad\|\pi(q)v-v\|<\eps\ (q\in Q).
\]
A nonzero invariant vector satisfies \(\pi(g)v=v\) for all \(g\in G\).

The unnumbered Question on PDF page 85 asks whether \(\operatorname{Mod}_S\) has property \((T)\). The source says a positive answer would answer a preceding question positively, but that question is absent and no specific implication is reconstructed. It describes Lie-group/representation-theoretic approaches to known discrete-group examples and suggests newer approaches of Shalom, Żuk and Valette, while reporting no serious work in that direction to the author's then-current knowledge. Its statement that the problem is completely open is a source-era assessment.

The supplied research summary instead claims a negative solution for closed genus-\(g\) groups with \(g\ge2\), without a primary citation. Its reasons, namely unbounded actions, vaguely invoked first Betti numbers and infinite abelian quotients of unspecified subgroups, do not establish that claim: property \((T)\) does not prohibit arbitrary unbounded actions or infinite abelian quotients of arbitrary subgroups. Passing to finite-index subgroups or to particular unitary actions would require additional stated hypotheses. Thus the summary's purported negative resolution and its changed surface range are preserved as conflicting, unverified status assertions, not adopted as the answer to the supplied question.
\subsection{Short English statement}
Does the source's mapping class group have Kazhdan's property \((T)\)? The source calls it open, while an uncited summary claims a negative solution using insufficiently specified reasoning. The draft preserves the exact property and records the status conflict without declaring a resolution.
\subsection{Sources}
\begin{itemize}
\item Local UnsolvedMath registry, permanent ID \(11000101\), exact source code \texttt{catalogue entry}, unnumbered Question, PDF page 85. No individual proposer or proposed year is supplied.
\item Farb, edited volume \emph{Problems on Mapping Class Groups and Related Topics}: \url{https://www.math.uchicago.edu/~farb/papers/mcgbook.pdf}. The volume has 329 explicitly typeset numbered/unnumbered Problem, Conjecture and Question headings, with page-specific identity used here.
\item Internal references: Y. Shalom [Shal], A. Żuk [Z], and A. Valette's N. Bourbaki Seminar report [V]. Full titles and bibliography entries are absent. The previous question invoked by the source is not included.
\item Extraction pdf-font-structure, accessed 2026-07-29. Source evidence NEEDS\_REVIEW and registry status open/research classification OPEN-TRIAGE require a current primary-literature audit. The summary labels itself SOLVED-IN-LITERATURE negatively for \(g\ge2\), but gives no verified theorem/reference and mentions unspecified Deligne-type or torsion-free finite-quotient subtleties. These do not supply the missing preceding question or settle this record.
\item Dataset attribution: Ulam AI / UnsolvedMath Contributors, CC BY 4.0. Source rights NEEDS\_REVIEW: publicly accessible author/publisher copy, redistribution terms requiring review.
\end{itemize}
\end{document}
